\documentclass[main.tex]{subfiles}


\begin{document}

% \AddToShipoutPictureFG{%
% \begin{tikzpicture}[overlay,remember picture]
% \path (current page.south west) -- (current page.north east)
% node[midway,scale=1,color=gray,sloped,opacity=0.4] {\textnormal{Кравченко Игорь Игоревич\hspace{1cm} +79010144910\hspace{1cm} super.antig@yandex.ru}};
% \end{tikzpicture}
% }


% %from pagenumber to up
% \phantomsection 
% \hypertarget{toc}{}

 

 

 
   

\section{Взаимодействие и сила}

\sbox{0}{%
    \begin{tikzpicture}[font=\small,            line cap=round,                             >={Stealth[inset=0pt,round,length=3mm, angle'=20]},vec/.style={>={Stealth[inset=0pt,round,length=3.5mm, angle'=20]}}]


\filldraw[lightgray,semitransparent] (5,0) rectangle (8,-0.3);
\draw (5,0) -- (8,0);




%%%%%%%%%%%%%%%%%%%%%%%%%%%%%%%%%%%%
%person 2
% PERSON
\def\W{5.2}     % ground width
\def\L{3.2}     % length
\def\T{0}    % plank thickness
\def\H{2.0}     % human height
  \def\x{6} % human position
  \coordinate (O) at (0,0);
\begin{scope}[yshift=0cm]  
  \draw[thick] (\x,\H) circle(0.3) coordinate (H) node[yshift=0.3cm,above] {$\text{Н}$};
  \draw[thick] (H)++(-90:0.3) coordinate (N) to[out=-92,in=92]++ (0,-0.40*\H) coordinate (P);
  \draw[thick] (N)++(-95:0.03) to[out=-105,in=90]++ (-0.02*\W,-0.4*\H);
  \draw[thick] (N)++(-85:0.03) to[out=-80,in=190]++ (0.13*\W,-0.1*\H) node (rh) {}; %right hand
  \draw[thick] (P) to[out=-92,in=82] (\x-0.02*\W,\T);
  \draw[thick] (P) to[out=-80,in=90] (\x+0.02*\W,\T);
\end{scope}

%%%%%%%%%%%%%%%%%%%%%%%%%%%%%%%%%%
%ball
\shade[ball color=lightgray] ([yshift=.3cm]rh) circle (.3) node[xshift=.3cm,right] {Ш};

    \end{tikzpicture}%
}

{%
\setlength\intextsep{0pt}
\begin{wrapfigure}[]{r}{\wd0}
    \centering
    \usebox{0}%
    \caption{Наблюдатель\\ и шар}
    \label{fig:p31}
\end{wrapfigure}

\textbf{Взаимодействие}~--- это взаимное действие двух тел друг на друга, приводящее к \emph{изменению их скоростей} (рис.\,\ref{fig:p31}).

% \textbf{Взаимодействие}~--- это влияние двух тел друг на друга, приводящее к \emph{изменению их скоростей} (рис.\,\ref{fig:p31}).

Из рис.\,\ref{fig:p31} видно, что рука наблюдателя~$\text{Н}$ действует на шар~Ш, иначе шар бы падал. С другой стороны, шар действует на руку, так как мышцы руки напряжены. Итак, в паре тел <<рука-шар>> происходит взаимодействие. (Однако изменения скоростей не наблюдаются из-за того, что каждое действие на данное тело компенсируется другим действием: действие руки на шар уравновешивается действием планеты, а действие шара на руку~--- действием мышц.)

}
% \begin{figure}[H]
%     \centering

%     \begin{tikzpicture}[font=\small,            line cap=round,                             >={Stealth[inset=0pt,round,length=3mm, angle'=20]},vec/.style={>={Stealth[inset=0pt,round,length=3.5mm, angle'=20]}}]


% \filldraw[lightgray,semitransparent] (5,0) rectangle (8,-0.3);
% \draw (5,0) -- (8,0);




% %%%%%%%%%%%%%%%%%%%%%%%%%%%%%%%%%%%%
% %person 2
% % PERSON
% \def\W{5.2}     % ground width
% \def\L{3.2}     % length
% \def\T{0}    % plank thickness
% \def\H{2.0}     % human height
%   \def\x{6} % human position
%   \coordinate (O) at (0,0);
% \begin{scope}[yshift=0cm]  
%   \draw[thick] (\x,\H) circle(0.3) coordinate (H) node[yshift=0.3cm,above] {$\text{Н}$};
%   \draw[thick] (H)++(-90:0.3) coordinate (N) to[out=-92,in=92]++ (0,-0.40*\H) coordinate (P);
%   \draw[thick] (N)++(-95:0.03) to[out=-105,in=90]++ (-0.02*\W,-0.4*\H);
%   \draw[thick] (N)++(-85:0.03) to[out=-80,in=190]++ (0.13*\W,-0.1*\H) node (rh) {}; %right hand
%   \draw[thick] (P) to[out=-92,in=82] (\x-0.02*\W,\T);
%   \draw[thick] (P) to[out=-80,in=90] (\x+0.02*\W,\T);
% \end{scope}

% %%%%%%%%%%%%%%%%%%%%%%%%%%%%%%%%%%
% %ball
% \shade[ball color=lightgray] ([yshift=.3cm]rh) circle (.3) node[xshift=.3cm,right] {Ш};


%     \end{tikzpicture}
% \caption{Наблюдатель и шар}
% \label{fig:p31}
% \end{figure}


Известны четыре типа взаимодействия в природе.
\begin{enumerate}
    \item \textbf{Гравитационное}: массивные тела взаимно притягивают друг друга.

    \emph{Пример:} планета и ее спутник.
    \item \textbf{Электромагнитное}: влияния между зарядами и/или магнитами.
    
    \emph{Пример:} два магнита.
    \item \textbf{Сильное}: влияния, <<удерживающие>> части атомного ядра.

    \emph{Пример:} протон и нейтрон.
    \item \textbf{Слабое}: влияния между любыми\footnote{Кроме фотонов.} элементарными частицами.  

    \emph{Пример:} нейтрон и электрон.
\end{enumerate}


% \sbox{0}{%
%     \begin{tikzpicture}[font=\small,            line cap=round,                             >={Stealth[inset=0pt,round,length=3mm, angle'=20]},vec/.style={>={Stealth[inset=0pt,round,length=3.5mm, angle'=20]}}]


% \filldraw[lightgray,semitransparent] (5,0) rectangle (8,-0.3);
% \draw (5,0) -- (8,0);

% \shade[ball color=lightgray,yshift=3cm] (6,.3) circle (.3);

% %vec
% \node at (6,3.3) [circle,fill=red,inner sep=0pt,minimum size=1mm,label=below:] {};
% \draw[->,red,thick,vec] (6,3.3) --++ (0,-1) node[black,right] {$\vec F_\text{планеты}$};

% \node at (6,3.3) [circle,fill=red,inner sep=0pt,minimum size=1mm,label=below:] {};
% \draw[->,red,thick,vec] (6,3.3) --++ (0,.75) node[black,right] {$\vec F_\text{воздуха}$};

% \draw[->,NavyBlue,thick,vec] (5.5,3.8) --++ (0,-1) node[black,left,midway] {$\vec v$};

%     \end{tikzpicture}%
% }

% {%
% \setlength\intextsep{0pt}
% \begin{wrapfigure}[]{r}{\wd0}
%     \centering
%     \usebox{0}%
%     \caption{Падение легкого шара}
%     \label{fig:p32}
% \end{wrapfigure}

\textbf{Сила} ($\vec F$ [Н])~--- это характеристика действия, показывающая, как велико воздействие на тело.

На рис.\,\ref{fig:p32} показан легкий шар в состоянии падения. 

\begin{figure}[H]
    \centering

    \begin{tikzpicture}[font=\small,            line cap=round,                             >={Stealth[inset=0pt,round,length=3mm, angle'=20]},vec/.style={>={Stealth[inset=0pt,round,length=3.5mm, angle'=20]}}]


\filldraw[lightgray,semitransparent] (5,0) rectangle (8,-0.3);
\draw (5,0) -- (8,0);

\shade[ball color=lightgray,yshift=1.5cm] (6,.3) circle (.3);

%vec
\node at (6,1.8) [circle,fill=red,inner sep=0pt,minimum size=1mm,label=below:] {};
\draw[->,red,thick,vec] (6,1.8) --++ (0,-1) node[black,right] {$\vec F_\text{планеты}$};

\node at (6,1.8) [circle,fill=red,inner sep=0pt,minimum size=1mm,label=below:] {};
\draw[->,red,thick,vec] (6,1.8) --++ (0,.75) node[black,right] {$\vec F_\text{воздуха}$};

\draw[->,NavyBlue,thick,vec] (5.5,2.3) --++ (0,-1) node[black,left,midway] {$\vec v$};

    \end{tikzpicture}%
\caption{Падение легкого шара}
\label{fig:p32}
\end{figure}

Силы, действующие на шар со стороны воздуха и планеты, обозначены красными векторами~$\vec F_\text{воздуха}$ и~$\vec F_\text{планеты}$ соответственно. 

\begin{law}
\textbf{Принцип суперпозиции сил.} Если на тело действуют силы~$\vec F_1, \vec F_2,\ldots,$ то их можно заменить одной силой 
\begin{equation}
\vec R\hm=\vec F_1+\vec F_2+\ldots
\end{equation}
\end{law}

Силу~$\vec R$ называют \textbf{результирующей} силой или \emph{равнодействующей} сил. Так, для шара на рис.\,\ref{fig:p32} результирующая сила равна: $\vec R_\text{ш}= \vec F_\text{воздуха}\hm+\vec F_\text{планеты}$. Из рис.\,\ref{fig:p32} также можно видеть, что $F_\text{воздуха}< F_\text{планеты}$, значит~--- эта <<суммарная>> сила~$\vec R_\text{ш}$
направлена вниз.















 
\end{document}